\documentclass[a4paper,11pt]{ltjsarticle}
\usepackage[margin=24mm,headheight=16pt]{geometry}
\usepackage{amsmath,amssymb,amsthm,mathtools,bm,mathrsfs}
\usepackage{booktabs,longtable,array,enumitem}
\usepackage{xcolor,needspace}
\usepackage[unicode,colorlinks=true,linkcolor=blue!40!black,citecolor=blue!40!black,urlcolor=blue!40!black]{hyperref}
\usepackage{fancyhdr}
\setlength{\parindent}{1em}
\setlength{\parskip}{0.2em}
\setlist{itemsep=0.2em,topsep=0.4em}
\allowdisplaybreaks[2]
\pagestyle{fancy}\fancyhf{}
\fancyhead[L]{\small 一楕円ブロックの非可換cubicと可換復元}
\fancyhead[R]{\small 2026-11-09 / v5}
\fancyfoot[C]{\thepage}
\newcommand{\CC}{\mathbb C}
\newcommand{\ZZ}{\mathbb Z}
\newcommand{\PP}{\mathbb P}
\newcommand{\OO}{\mathcal O}
\newcommand{\cK}{\mathcal K}
\newcommand{\Perf}{\operatorname{Perf}}
\newcommand{\Ku}{\operatorname{Ku}}
\newcommand{\MF}{\operatorname{MF}}
\newcommand{\HH}{\operatorname{HH}}
\newcommand{\qgr}{\operatorname{qgr}}
\newcommand{\Dsg}{D_{\mathrm{sg}}}
\newcommand{\RHom}{\operatorname{RHom}}
\newcommand{\Hom}{\operatorname{Hom}}
\newcommand{\Ext}{\operatorname{Ext}}
\newcommand{\Spec}{\operatorname{Spec}}
\newcommand{\Id}{\operatorname{Id}}
\newcommand{\sym}{\operatorname{sym}}
\newcommand{\Cone}{\operatorname{Cone}}
\newcommand{\thick}{\operatorname{thick}}
\newcommand{\Pic}{\operatorname{Pic}}
\newcommand{\otL}{\mathbin{\otimes^{\mathbf L}}}
\theoremstyle{plain}
\newtheorem{theorem}{定理}[section]
\newtheorem{proposition}[theorem]{命題}
\newtheorem{lemma}[theorem]{補題}
\newtheorem{corollary}[theorem]{系}
\theoremstyle{definition}
\newtheorem{definition}[theorem]{定義}
\newtheorem{example}[theorem]{例}
\theoremstyle{remark}
\newtheorem{remark}[theorem]{注意}
\renewcommand{\proofname}{証明}
\hypersetup{pdftitle={一楕円ブロックの可換復元：研究ノート}}
\title{研究ノート\\\large 一楕円ブロックの可換復元と相対化の課題}
\author{}
\date{2026年11月9日・v5}
\begin{document}
\maketitle

\section{現在の知識と主張の状態}
本ノートの正本はこのTeX、詳細な証明の正本は
\texttt{elliptic\_cubic\_reconstruction.pdf}である。
本一式は\texttt{nc-cubic}の一楕円ブロック部分を独立に再開するための自己完結した部分問題bundleであり、
一般Poisson分類・退化スキーム・多ブロック生成などを含む親プロジェクト全体の置換版ではない。
親リポジトリの既存の有効な結果を削除しない。

\begin{longtable}{p{0.27\linewidth}p{0.65\linewidth}}
\toprule
対象 & 現在の状態・証明の所在\\\midrule
一般Hesse係数の一ブロック & 論文定理1.1。中心部分$P(u)$が分離され孤立臨界点を持てば、
各良い複素ファイバーでrotation付き可換模型を構成。平行移動の有限位数性は不要。\\
任意の滑らかな三変数三次式 & 論文「主定理の証明」の系。Hesse変換と
$h'=\lambda\det(M)h$による。中心線形シフトで生じる混合項も含む。\\
復元cubicの変形 & 論文第5節。$h$では一次が零、二次が非零。
$q=h^2$でのJacobian類を明示。coarse moduliにおける最初の次数とは区別。\\
一定の可換模型を持つ代数的曲線 & 論文第6節。$\vartheta(h,t)=\tau$の曲線で、
非可換代数の一次変形は非自明だが、各ファイバーの残余pairは同じ可換cubicのpair。
相対自明化は主張しない。\\
中心商の相対表示 & 論文第7節。$S_R/(\Omega_R)$の相対THCR表示と任意の底変更との両立。
残余圏の相対saturationを証明したものではない。\\
二つ以上の無限位数ブロック & 全生成の一般的証明は未完成。今回の一ブロック定理から無条件には従わない。\\
一般の本質的な中心混合項 & 未解決。中心変数の通常の偏微分が非可換係数代数の導分にならない場合がある。\\
\bottomrule
\end{longtable}

ここで「証明を記載」「主要論証を今回点検」「厳密な有限計算を実行」は別の状態である。
本稿は公刊・査読済みではなく、証明支援系による形式化も未実施である。
外部定理の適用範囲と新たな組合せを区別し、新規性の初出確認は継続中とする。

\section{既存原稿との論理的な関係}
確認した親リポジトリの基点コミットは
\texttt{e61c790f0101585522fa18d8c353d70d2d9bb8c6}である。
同版の\texttt{tex/papers/rotation\_models.tex}には
\texttt{cor:elliptic-central-cubic}として一般の楕円ブロックの比較が既にある。
一方、Hesse中心元の節では一般$t$の可換模型が候補のままになっていた。
本稿はこの既存の一般比較の仮定を明示的Hesse族で検証して接続したものであり、
新しい抽象的な比較定理を初めて発見したという位置づけではない。

その接続を旧原稿への参照だけで代用せず、中心商、rotationの共役、外部積の生成、
有限次数差の降下を本稿に記した。
二次変形の非自明性、一定模型曲線、相対THCR表示は、それぞれ独立した証明を追加している。
親プロジェクトの一般AS-regular中心三次商に対するZhang twist反例は消えていない。
今回の存在定理は指定した楕円ブロック族に限られ、それらの反例とは適用範囲が異なる。

\section{符号、正規化、coherenceの注意}
\subsection{中心元と点曲線}
$S_{h,t}$のSklyaninパラメータは
$(\alpha,\beta,\gamma)=(1-ht/2,-1-ht/2,-h)$である。
二次関係式の多重線形化を$o=[1:-1:0]$で評価すると、核は
$p=[\alpha:\beta:\gamma]$になる。
これは右加群の規約で$\sigma=t_p$を用いることと整合する。
$h=0$の周囲の代数のpoint schemeは$\PP^2$だが、その中心商は$E_t$である。

中心三次元についてはDe Laet--Le BruynのarXiv v1のSection 2, p.~3の対称表示を用いる。
同版冒頭のTheorem 1の表示にはSection 2との係数差がある。
このため、古典的中心性を引用するだけでなく、
$G=-9h\Omega_{h,t}$という変換恒等式を別途厳密計算した。
新規な普遍的中心元の公式を主張するのではなく、既知の式の量子化パラメータへの翻訳である。

\subsection{rotationと周期性}
右加群のtailsのシフトは$\sigma_*(-\otimes L)$。
windowで$T_{\OO_E}$を加えると
$\rho_p=t_{-p}^*\rho_0$、$\rho_0=T_{\OO_E}(-\otimes L)$となる。
次数零直線束$P_q$によるテンソル$A_q$は
$A_q\rho_0A_q^{-1}=t_q^*\rho_0$を満たすので、
$A_p$が非可換側から可換側へのintertwinerになる。
単に$q=p$か$-p$かを記憶するのでなく、この向きで確認する。

$(3)\simeq[2]$の二つの規格化の差はスカラーである。
その三乗根を選ぶことで有限群降下のcoherenceを合わせる。
この選択は各複素ファイバーでは可能だが、底全体上で代数的に一意に選べるという意味ではない。
相対問題では可逆層やétale被覆に沿う選択の問題を別途処理する必要がある。

\subsection{粗化次数の対象を落とさない}
中心変数を一因子にまとめた次数群は
$\Gamma=\ZZ^2/(3e_0-3e_1)$であり、その通常次数への核は有限である。
一方、多ブロックで元の細分次数から$\Gamma$へ移す核は無限格子を含む。
多重tailsが可換な積の幾何を与えることだけから、粗化次数付き全加群の生成性を結論してはならない。
本稿の一ブロックでは、中心Jacobian支持へのKoszul帰着と
$\Gamma/\ZZ e_0\cong\ZZ/3$の分解を使って全対象を扱う。

\section{相対復元の次の段階}
\subsection{現在すでにある代数的入力}
$R=\CC[h,(Ub\Delta)^{-1}]$上に、次数別自由な$S_R$、中心正則元$\Omega_R$、
滑らかな楕円族$\mathscr E$、次数三の線束$\mathcal L$、平行移動$\sigma$が存在する。
論文第7節は$S_R/(\Omega_R)$の相対THCR表示と底変更を与える。
従って次の相対windowの議論に必要な中心商は、形式完備化ではなく実際の有限型底上で得られている。

\subsection{具体的な残課題}
最初に$S_R$とその中心商の左右noetherian性、および必要な局所cohomologyの有限性を示す。
その後、$A_0=R$を持つ次数付きGorenstein代数の意味で
左右の有限injective dimensionと
$\RHom_A(A_0,A)\simeq A_0[n](a)$を検証する。
Burke--StevensonのDefinition 5.1とTheorem 6.4はこの比較に適した出発点である。
ただし同定理の代数の条件を、各閉ファイバーのAS-regular性と同一視しない。

次に比較をrelative perfect対象のdg模型で構成し、
Hom、射影関手、rotation、三次周期性が導来底変更と両立することを示す。
中心$P(u)$を加えるKoszul生成の証明を底上で行い、
最後に有限群の同変化と底変更を両立させる。
ここまで完成すれば、非可換の相対残余圏を既に存在するsmooth properな
$\Ku(Y/R)$と同定することで相対saturationを得る見込みがある。

別の経路として、親原稿の全超曲面圏の相対proper模型に対角の有限perfect証明書を作り、
有限相対例外列を除く前後のsmoothnessの同値を利用できる。
本稿はどちらの経路も全て完了したとはしていない。

\section{計算証明書と再現可能性}
\texttt{computations/verify\_reconstruction.py}はSymPyによる有理係数の恒等式計算と、
$\mathbb F_{1009}$での線形代数を行う。
確認する事項は中心式の係数変換、Sklyanin判別式、点の所属と作用の向き、
$\vartheta-t$、一定模型曲線の式とétale性、Fermat特殊化である。
さらに、$x_i\partial_{x_j}f_t$の三次係数行列から$xyz$列を除いた$9\times9$小行列式は
\[
 -(1+t^3)^2
\]
となる。これは滑らかな全Hesse係数で$[xyz]\ne0$であることの有限証明書である。
有限体検算は8個の特殊化で、中心元と三生成元との交換子が次数四の二次関係式イデアルに属することを確認する。
これらは圏同値・全生成性・相対化を形式検証したものではない。

\end{document}